% =============================================================================
% TP 1.2 - Estudio Económico-Matemático de Apuestas en la Ruleta
% Universidad Tecnológica Nacional – FRRO
% Materia: Simulación 2026
%
% Compilación (desde la carpeta tp1.2/informe/):
%   pdflatex main.tex
%   bibtex main
%   pdflatex main.tex
%   pdflatex main.tex
%
% IMPORTANTE: antes de compilar, generar los gráficos ejecutando desde la raíz:
%   python tp1.2/ruleta1_2.py -c 5 -n 500 -s m -a i -ci 100
%   python tp1.2/ruleta1_2.py -c 5 -n 500 -s m -a f -ci 100
%   python tp1.2/ruleta1_2.py -c 5 -n 500 -s d -a i -ci 100
%   python tp1.2/ruleta1_2.py -c 5 -n 500 -s d -a f -ci 100
%   python tp1.2/ruleta1_2.py -c 5 -n 500 -s f -a i -ci 100
%   python tp1.2/ruleta1_2.py -c 5 -n 500 -s f -a f -ci 100
%   python tp1.2/ruleta1_2.py -c 5 -n 500 -s o -a i -ci 100
%   python tp1.2/ruleta1_2.py -c 5 -n 500 -s o -a f -ci 100
%
% En Overleaf: subir main.tex, referencias.bib y la carpeta graficos/ completa.
% Aplicar el template Cornell/arXiv (pkzcrhzcdxmc).
% =============================================================================

\documentclass[a4paper, 11pt]{article}

% --- Codificación y lenguaje ---
\usepackage[utf8]{inputenc}
\usepackage[T1]{fontenc}
\usepackage[spanish, es-nodecimaldot]{babel}

% --- Matemática ---
\usepackage{amsmath}
\usepackage{amssymb}

% --- Gráficos y figuras ---
\usepackage{graphicx}
\usepackage{float}
\usepackage{caption}
\usepackage{subcaption}

% --- Tablas ---
\usepackage{booktabs}
\usepackage{array}
\usepackage{multirow}

% --- Formato ---
\usepackage[margin=2.5cm]{geometry}
\usepackage{parskip}
\usepackage{microtype}
\usepackage{lmodern}

% --- Hipervínculos ---
\usepackage[
  colorlinks=true,
  linkcolor=blue!70!black,
  citecolor=blue!70!black,
  urlcolor=blue!70!black
]{hyperref}

% --- Código fuente ---
\usepackage{listings}
\usepackage{xcolor}

\lstdefinestyle{python}{
  language=Python,
  basicstyle=\ttfamily\small,
  keywordstyle=\color{blue!80!black}\bfseries,
  commentstyle=\color{gray},
  stringstyle=\color{orange!80!black},
  numbers=left,
  numberstyle=\tiny\color{gray},
  numbersep=8pt,
  breaklines=true,
  frame=single,
  backgroundcolor=\color{gray!8},
  showstringspaces=false,
  tabsize=4,
  captionpos=b
}

% =============================================================================

\title{%
  \textbf{TP 1.2 -- Estudio Económico-Matemático de Apuestas en la Ruleta} \\[6pt]
  \large Universidad Tecnológica Nacional -- FRRO \\
  \large Materia: Simulación -- Año 2026
}

\author{
  Ceballos, R. \\
  \texttt{rceballos.work@gmail.com}
}

\date{13 de mayo, 2026}

% =============================================================================

\begin{document}

\maketitle

% ---------------------------------------------------------------------------
\begin{abstract}
El presente trabajo extiende la simulación de la ruleta europea desarrollada
en el TP~1.1 para analizar cuatro estrategias de apuesta: Martingala,
D'Alembert, Fibonacci y Paroli (estrategia inversa, de elaboración propia).
Para cada estrategia se simula el flujo de caja bajo dos supuestos mutuamente
excluyentes: capital infinito (ideal) y capital finito (real), este último
con registro de bancarrotas. Las métricas analizadas son la frecuencia relativa
acumulada de apuestas favorables (\textit{frsa}) y la evolución del capital
disponible (\textit{cc}) a lo largo de $n = 500$ tiradas por corrida, con
$c = 5$ corridas independientes por estrategia y modo de capital.
Los resultados evidencian que, independientemente de la estrategia empleada,
la ventaja de la casa del $1/37 \approx 2{,}7\,\%$ genera una pérdida
esperada de capital que ninguna estrategia de apuesta puede revertir a
largo plazo. La principal diferencia entre estrategias radica en el perfil
de riesgo y la velocidad de descenso del capital.
\end{abstract}

% ---------------------------------------------------------------------------
\section{Introducción}

El estudio de estrategias de apuesta en juegos de azar ha sido objeto de
análisis matemáticos y probabilísticos desde el siglo XVIII. La ruleta,
en particular, ha inspirado numerosos sistemas de apuesta que prometen
maximizar las ganancias o minimizar las pérdidas. Sin embargo, desde la
perspectiva de la teoría de la probabilidad, cualquier apuesta en una
ruleta europea tiene la misma ganancia esperada negativa, independientemente
de la estrategia empleada \cite{ross2014introduction}.

La \textit{ventaja de la casa} (\textit{house edge}) en la ruleta europea
es del $1/37 \approx 2{,}7\,\%$: por cada unidad apostada, el jugador
pierde en promedio $1/37$ unidades a largo plazo, sin excepción
\cite{wikipedia_ruleta}. Esta propiedad es una consecuencia directa de la
linealidad del valor esperado: si $E[\text{ganancia por tirada}] < 0$,
ninguna función de las apuestas pasadas puede invertir este resultado.

El propósito de este trabajo es cuantificar y visualizar el comportamiento
de cuatro estrategias de apuesta (Martingala, D'Alembert, Fibonacci y Paroli)
mediante simulación computacional, verificando la convergencia de las
métricas empíricas hacia los valores teóricos y comparando los perfiles de
riesgo de cada estrategia bajo condiciones de capital infinito y finito
\cite{sportium_estrategias}.

% ---------------------------------------------------------------------------
\section{Métodos}

\subsection{Modelo de la ruleta y tipos de apuesta}

Se utiliza el mismo modelo de ruleta europea del TP~1.1: espacio muestral
$\Omega = \{0, 1, 2, \ldots, 36\}$, con distribución uniforme discreta:

\begin{equation}
  P(X = k) = \frac{1}{37}, \quad k \in \Omega
  \label{eq:prob_uniforme}
\end{equation}

En este trabajo se considera la apuesta a \textbf{color rojo} (cuota 1:1),
que cubre 18 números: $\{1, 3, 5, 7, 9, 12, 14, 16, 18, 19, 21, 23, 25,
27, 30, 32, 34, 36\}$. La probabilidad de ganar esta apuesta es:

\begin{equation}
  P(\text{victoria}) = \frac{18}{37} \approx 0{,}4865
  \label{eq:prob_color}
\end{equation}

La ganancia esperada por unidad apostada en cada tirada es:

\begin{equation}
  E[\text{ganancia}] = \frac{18}{37} \cdot 1 + \frac{19}{37} \cdot (-1)
    = -\frac{1}{37} \approx -0{,}027
  \label{eq:ganancia_esperada}
\end{equation}

Esto implica que, tras $n$ tiradas apostando $a_i$ en la tirada $i$, el
capital esperado decrece según:

\begin{equation}
  E[cc_n] = cc_0 - \frac{1}{37} \sum_{i=1}^{n} a_i
  \label{eq:capital_esperado}
\end{equation}

donde $cc_0$ es el capital inicial. Cuanto mayor sea el monto apostado (como
en Martingala), más rápida será la pérdida esperada.

\subsection{Estrategias de apuesta}

Todas las estrategias parten de una apuesta base $a_{\text{base}} = 1$ unidad
y determinan la próxima apuesta $a_{n+1}$ en función del resultado de la
tirada $n$.

\subsubsection{Martingala}
\label{sec:martingala}

La Martingala es la estrategia de progresión negativa más conocida: ante
cada derrota se duplica la apuesta, y ante una victoria se reinicia al
valor base \cite{wikipedia_martingala}:

\begin{equation}
  a_{n+1} =
  \begin{cases}
    a_{\text{base}} & \text{si la tirada } n \text{ fue victoria} \\
    2 \cdot a_n      & \text{si la tirada } n \text{ fue derrota}
  \end{cases}
  \label{eq:martingala}
\end{equation}

Tras $k$ derrotas consecutivas la apuesta es $a_k = 2^k \cdot a_{\text{base}}$.
Si la siguiente tirada resulta victoria, la ganancia neta acumulada del ciclo
siempre es $+1$ unidad:

\begin{equation}
  \text{Ganancia neta} = 2^k - \sum_{j=0}^{k-1} 2^j = 2^k - (2^k - 1) = 1
  \label{eq:martingala_ganancia}
\end{equation}

El riesgo es el crecimiento \textbf{exponencial} del capital requerido.

\subsubsection{D'Alembert}
\label{sec:dalembert}

La D'Alembert es una progresión aritmética, más conservadora que la
Martingala \cite{wikipedia_dalembert}:

\begin{equation}
  a_{n+1} =
  \begin{cases}
    \max(a_{\text{base}},\; a_n - 1) & \text{si victoria} \\
    a_n + 1                           & \text{si derrota}
  \end{cases}
  \label{eq:dalembert}
\end{equation}

Tras $k$ derrotas consecutivas: $a_k = a_{\text{base}} + k$ (crecimiento lineal).
El sistema se equilibra cuando el número de victorias iguala al de derrotas
desde el inicio del ciclo.

\subsubsection{Fibonacci}
\label{sec:fibonacci}

La estrategia de Fibonacci recorre la secuencia
$\mathcal{F} = (1, 1, 2, 3, 5, 8, 13, 21, 34, 55, \ldots)$
definida por $F_1 = F_2 = 1$ y $F_n = F_{n-1} + F_{n-2}$
\cite{wikipedia_fibonacci_bet}:

\begin{equation}
  \text{índice}_{n+1} =
  \begin{cases}
    \max(0,\; \text{índice}_n - 2) & \text{si victoria} \\
    \text{índice}_n + 1             & \text{si derrota}
  \end{cases}
  \quad \Longrightarrow \quad a_{n+1} = F_{\,\text{índice}_{n+1}}
  \label{eq:fibonacci}
\end{equation}

El retroceso de dos posiciones en la victoria permite recuperar las pérdidas
acumuladas con menor crecimiento de la apuesta que en la Martingala. El
crecimiento sigue una progresión sub-exponencial (la razón áurea $\varphi
\approx 1{,}618$).

\subsubsection{Paroli (estrategia propia)}
\label{sec:paroli}

La Paroli es la estrategia inversa de la Martingala: se duplica la apuesta
en las victorias (hasta un máximo de 3 victorias consecutivas) y se reinicia
en cualquier derrota o al completar el ciclo de 3 victorias:

\begin{equation}
  a_{n+1} =
  \begin{cases}
    2 \cdot a_n        & \text{si victoria y victorias consecutivas} < 3 \\
    a_{\text{base}}    & \text{si victoria y victorias consecutivas} \geq 3 \\
    a_{\text{base}}    & \text{si derrota}
  \end{cases}
  \label{eq:paroli}
\end{equation}

La secuencia de apuestas por ciclo es $1 \to 2 \to 4 \to 1$. El resultado
de cada ciclo es:

\begin{itemize}
  \item \textbf{Derrota en cualquier tirada del ciclo:} pérdida de
        $-a_{\text{base}} = -1$ unidad (siempre).
  \item \textbf{Victoria en las 3 tiradas:} ganancia neta de
        $+1 + 2 + 4 = +7$ unidades.
\end{itemize}

La probabilidad de completar el ciclo con 3 victorias seguidas es:

\begin{equation}
  P(\text{ciclo exitoso}) = \left(\frac{18}{37}\right)^3 \approx 0{,}1136
  \label{eq:paroli_prob}
\end{equation}

La ganancia esperada por ciclo sigue siendo negativa:

\begin{equation}
  E[\text{ganancia ciclo}] = 7 \cdot \left(\frac{18}{37}\right)^3
    - 1 \cdot \left[1 - \left(\frac{18}{37}\right)^3\right] \approx -0{,}09
  \label{eq:paroli_ev}
\end{equation}

La ventaja de esta estrategia no es en valor esperado, sino en el
\textbf{control del riesgo}: la pérdida máxima por tirada individual
es siempre de 1 unidad base, mientras que la potencial ganancia
(en ciclo exitoso) es de 7 unidades.

\subsection{Modos de capital}

\paragraph{Capital infinito (ideal).}
El jugador no tiene límite de fondos. El capital puede volverse negativo
sin restricciones. Este supuesto permite observar el comportamiento puro
de la estrategia sin interrupciones.

\paragraph{Capital finito (real).}
El jugador parte de un capital inicial $cc_0$. Si en alguna tirada el
capital disponible es menor que la próxima apuesta requerida por la
estrategia, se produce una \textbf{bancarrota}: se registra el evento,
el capital se reinicia a $cc_0$ y la estrategia vuelve a su apuesta base.
El número total de bancarrotas por corrida es un indicador del riesgo
práctico de la estrategia.

\subsection{Métricas de análisis}

\paragraph{Frecuencia relativa acumulada de apuesta favorable (frsa).}

\begin{equation}
  frsa_i = \frac{\#\{j \leq i : \text{tirada } j \text{ fue victoria}\}}{i}
  \label{eq:frsa}
\end{equation}

Por la Ley de los Grandes Números, $frsa_i \xrightarrow{p} 18/37 \approx 0{,}4865$
cuando $i \to \infty$.

\paragraph{Flujo de caja (cc).}

Evolución del capital disponible tirada a tirada:

\begin{equation}
  cc_i =
  \begin{cases}
    cc_{i-1} + a_i & \text{si la tirada } i \text{ fue victoria} \\
    cc_{i-1} - a_i & \text{si la tirada } i \text{ fue derrota}
  \end{cases}
  \label{eq:flujo_caja}
\end{equation}

La línea de referencia es $cc_0$ (capital inicial, \textit{fci}): si
$cc_i > cc_0$ el jugador está en ganancia; si $cc_i < cc_0$, en pérdida.

\subsection{Parámetros de la simulación}

\begin{itemize}
  \item Corridas por estrategia y modo: $c = 5$.
  \item Tiradas por corrida: $n = 500$.
  \item Capital inicial (modo finito): $cc_0 = 100$ unidades.
  \item Apuesta base: $a_{\text{base}} = 1$ unidad.
  \item Tipo de apuesta: color rojo (cuota 1:1).
\end{itemize}

\subsection{Implementación}

El programa se implementa en Python~3 en el archivo
\texttt{tp1.2/ruleta1\_2.py}. Cada estrategia es una clase independiente
con el método \texttt{siguiente(gano)} que devuelve la próxima apuesta.
A continuación se muestra la implementación de la Martingala y el núcleo
del ciclo de simulación como representativos del diseño general:

\begin{lstlisting}[style=python, caption={Clase Martingala y ciclo principal de simulación.}]
class Martingala:
    def __init__(self) -> None:
        self._apuesta: int = APUESTA_BASE

    def siguiente(self, gano: bool) -> int:
        self._apuesta = APUESTA_BASE if gano else self._apuesta * 2
        return self._apuesta

    def reset(self) -> None:
        self._apuesta = APUESTA_BASE


def simular_corrida(n_tiradas, codigo_estrategia,
                    capital_inicial, capital_infinito,
                    numero_elegido=None):
    estrategia = crear_estrategia(codigo_estrategia)
    capital = capital_inicial
    aciertos = 0
    bancarrotas = 0
    apuesta = APUESTA_BASE

    for i in range(n_tiradas):
        # Bancarrota: capital insuficiente para la apuesta
        if not capital_infinito and capital < apuesta:
            bancarrotas += 1
            capital = capital_inicial
            estrategia.reset()
            apuesta = APUESTA_BASE

        numero = girar_ruleta()
        gano, cuota = evaluar_apuesta(numero, numero_elegido)

        capital += apuesta * cuota if gano else -apuesta
        frsa.append(aciertos / (i + 1))
        apuesta = estrategia.siguiente(gano)

    return ResultadoCorrida(frsa, flujo_caja, capital_inicial, bancarrotas)
\end{lstlisting}

Los números aleatorios se generan con \texttt{random.randint(0, 36)}
del módulo estándar de Python (Mersenne Twister) \cite{python_random}.
Las visualizaciones se producen con Matplotlib \cite{matplotlib2007}.

% ---------------------------------------------------------------------------
\section{Resultados}

Para cada estrategia se presentan dos figuras: una corrida individual
representativa (corrida~1) y la superposición de las cinco corridas
independientes. Las figuras de capital infinito permiten observar la
tendencia estructural de la estrategia; las de capital finito ilustran
el impacto del límite de fondos y la ocurrencia de bancarrotas.

% ····················
\subsection{Martingala}

\begin{figure}[H]
  \centering
  \includegraphics[width=\linewidth]{../graficos/m_inf/corrida_01.png}
  \caption{Martingala -- capital infinito: corrida 1 ($n = 500$ tiradas).
    Izquierda: \textit{frsa} acumulada vs.\ valor esperado ($18/37$).
    Derecha: flujo de caja vs.\ capital inicial (\textit{fci}).}
  \label{fig:m_inf_c1}
\end{figure}

\begin{figure}[H]
  \centering
  \includegraphics[width=\linewidth]{../graficos/m_fin/corrida_01.png}
  \caption{Martingala -- capital finito ($cc_0 = 100$): corrida 1 ($n = 500$ tiradas).
    Los picos en el flujo de caja corresponden a victorias que recuperan
    rachas de derrotas. Las caídas abruptas indican bancarrotas seguidas
    de reinicio de capital.}
  \label{fig:m_fin_c1}
\end{figure}

\begin{figure}[H]
  \centering
  \includegraphics[width=\linewidth]{../graficos/m_fin/todas_las_corridas.png}
  \caption{Martingala -- capital finito: cinco corridas superpuestas.
    La dispersión en el flujo de caja es la mayor de todas las estrategias
    analizadas, dado el carácter exponencial de la progresión de apuestas.}
  \label{fig:m_fin_todas}
\end{figure}

% ····················
\subsection{D'Alembert}

\begin{figure}[H]
  \centering
  \includegraphics[width=\linewidth]{../graficos/d_inf/corrida_01.png}
  \caption{D'Alembert -- capital infinito: corrida 1 ($n = 500$ tiradas).
    La progresión aritmética de las apuestas produce un flujo de caja con
    menor volatilidad que la Martingala.}
  \label{fig:d_inf_c1}
\end{figure}

\begin{figure}[H]
  \centering
  \includegraphics[width=\linewidth]{../graficos/d_fin/todas_las_corridas.png}
  \caption{D'Alembert -- capital finito: cinco corridas superpuestas.
    La menor amplitud de las apuestas reduce la frecuencia de bancarrotas
    respecto a la Martingala con el mismo capital inicial.}
  \label{fig:d_fin_todas}
\end{figure}

% ····················
\subsection{Fibonacci}

\begin{figure}[H]
  \centering
  \includegraphics[width=\linewidth]{../graficos/f_inf/corrida_01.png}
  \caption{Fibonacci -- capital infinito: corrida 1 ($n = 500$ tiradas).
    La secuencia de Fibonacci crece más lentamente que la Martingala,
    generando un perfil de riesgo intermedio.}
  \label{fig:f_inf_c1}
\end{figure}

\begin{figure}[H]
  \centering
  \includegraphics[width=\linewidth]{../graficos/f_fin/todas_las_corridas.png}
  \caption{Fibonacci -- capital finito: cinco corridas superpuestas.
    Se observa una volatilidad intermedia entre Martingala y D'Alembert.}
  \label{fig:f_fin_todas}
\end{figure}

% ····················
\subsection{Paroli (estrategia inversa)}

\begin{figure}[H]
  \centering
  \includegraphics[width=\linewidth]{../graficos/o_inf/corrida_01.png}
  \caption{Paroli -- capital infinito: corrida 1 ($n = 500$ tiradas).
    El flujo de caja presenta saltos discretos de $+7$ unidades (ciclos exitosos)
    y caídas de $-1$ unidad (derrotas), generando el patrón de escalera
    característica de esta estrategia.}
  \label{fig:o_inf_c1}
\end{figure}

\begin{figure}[H]
  \centering
  \includegraphics[width=\linewidth]{../graficos/o_fin/todas_las_corridas.png}
  \caption{Paroli -- capital finito: cinco corridas superpuestas. Al apostar
    como máximo 4 unidades por tirada, el número de bancarrotas es el menor
    de las cuatro estrategias con el mismo capital inicial.}
  \label{fig:o_fin_todas}
\end{figure}

\subsection{Comparación de bancarrotas (capital finito)}

La Tabla~\ref{tab:bancarrotas} resume el número total de bancarrotas
registradas en las 5 corridas ($n = 500$ tiradas cada una, $cc_0 = 100$
unidades) para cada estrategia.

\begin{table}[H]
  \centering
  \caption{Total de bancarrotas en 5 corridas ($n = 500$, $cc_0 = 100$).
    Los valores son representativos de una ejecución; las ejecuciones
    posteriores variarán por el carácter estocástico del modelo.}
  \label{tab:bancarrotas}
  \begin{tabular}{lcc}
    \toprule
    \textbf{Estrategia} & \textbf{Apuesta máx.\ posible} & \textbf{Bancarrotas totales} \\
    \midrule
    Martingala   & Exponencial ($2^k$) & Ver gráfico \\
    D'Alembert   & Lineal ($k + 1$)    & Ver gráfico \\
    Fibonacci    & Sub-exponencial ($F_k$) & Ver gráfico \\
    Paroli       & Acotada ($\leq 4$)  & Ver gráfico \\
    \bottomrule
  \end{tabular}
\end{table}

% ---------------------------------------------------------------------------
\section{Discusión}

\subsection{Invariancia del valor esperado}

El resultado más importante de este análisis es que \textbf{ninguna
estrategia de apuesta puede modificar el valor esperado de ganancia},
que siempre es $-1/37$ por unidad apostada
\cite{walpole2012probabilidad, devore2016probabilidad}.

La demostración es directa: sea $a_i$ la apuesta en la tirada $i$,
determinada por cualquier función de los resultados anteriores
$(x_1, x_2, \ldots, x_{i-1})$. Por la linealidad del valor esperado
y la independencia de los giros:

\begin{equation}
  E\!\left[\sum_{i=1}^{n} g_i\right]
    = \sum_{i=1}^{n} E[g_i]
    = \sum_{i=1}^{n} E[a_i] \cdot \left(-\frac{1}{37}\right)
    = -\frac{1}{37} \sum_{i=1}^{n} E[a_i]
  \label{eq:ev_invarianza}
\end{equation}

donde $g_i$ es la ganancia (positiva o negativa) en la tirada $i$.
Como $E[a_i] > 0$ para toda estrategia, el valor esperado total siempre
es negativo. Las estrategias difieren en la \textit{varianza} de $a_i$,
no en su media.

\subsection{Perfiles de riesgo}

Las cuatro estrategias presentan perfiles de riesgo distintos, resumidos
en la Tabla~\ref{tab:riesgo}.

\begin{table}[H]
  \centering
  \caption{Comparación de perfiles de riesgo de las cuatro estrategias
    analizadas.}
  \label{tab:riesgo}
  \begin{tabular}{lcccc}
    \toprule
    \textbf{Estrategia} & \textbf{Crecim.\ apuesta}
      & \textbf{Pérd.\ máx.\ tirada}
      & \textbf{Gan.\ por victoria}
      & \textbf{Riesgo bancarrota} \\
    \midrule
    Martingala   & Exponencial & $2^k$ & $+1$ siempre & Muy alto \\
    D'Alembert   & Aritmético  & $k+1$ & Variable     & Moderado \\
    Fibonacci    & Sub-exp.    & $F_k$ & Variable     & Moderado \\
    Paroli       & Acotado ($\leq 4$) & $-1$ siempre & $+7$ (ciclo) & Muy bajo \\
    \bottomrule
  \end{tabular}
\end{table}

\subsection{Análisis de la Martingala}

La Martingala es matemáticamente ``segura'' en capital infinito: garantiza
siempre una ganancia de $+1$ unidad por ciclo, independientemente de cuántas
derrotas consecutivas ocurran. Sin embargo, bajo capital finito, la
probabilidad de una racha de $k$ derrotas consecutivas es:

\begin{equation}
  P(\text{k derrotas seguidas}) = \left(\frac{19}{37}\right)^k
  \label{eq:racha_derrotas}
\end{equation}

Con $cc_0 = 100$ y $a_{\text{base}} = 1$, el capital se agota si se
producen $k \geq 7$ derrotas consecutivas (pues $2^7 - 1 = 127 > 100$).
La probabilidad de esta racha es $(19/37)^7 \approx 0{,}79\,\%$ por ciclo,
lo que a lo largo de 500 tiradas se vuelve un evento frecuente.

\subsection{Análisis de la Paroli}

La Paroli invierte el paradigma de riesgo: la pérdida máxima por tirada
individual es siempre $-a_{\text{base}}$, lo que hace prácticamente
imposible la bancarrota con capitales moderados. La contrapartida es que
el ciclo exitoso (3 victorias consecutivas) tiene probabilidad solo del
$11{,}36\,\%$, lo que genera una larga secuencia de pérdidas unitarias
interrumpida por ganancias de $+7$ unidades. La ganancia esperada por
ciclo (ecuación~\ref{eq:paroli_ev}) sigue siendo negativa, pero el perfil
temporal es radicalmente diferente al de las progresiones negativas.

\subsection{Convergencia de la frsa}

Para todas las estrategias, la \textit{frsa} converge hacia $18/37 \approx
0{,}4865$ conforme $n$ crece, en concordancia con la Ley de los Grandes
Números. La velocidad de convergencia es idéntica para las cuatro estrategias
porque la probabilidad de victoria en cada tirada ($18/37$) no depende de
la apuesta realizada: la estrategia determina \textit{cuánto} se apuesta,
no \textit{con qué frecuencia} se gana.

% ---------------------------------------------------------------------------
\section{Conclusiones}

El análisis de las cuatro estrategias de apuesta en la ruleta europea
permite extraer las siguientes conclusiones:

\begin{enumerate}
  \item \textbf{Invariancia del resultado esperado:} todas las estrategias
        analizadas producen la misma ganancia esperada negativa de
        $-1/37 \approx -2{,}7\,\%$ por unidad apostada. Ningún sistema
        de apuesta puede superar la ventaja estructural de la casa en una
        ruleta justa.

  \item \textbf{La elección de estrategia es una elección de perfil de riesgo:}
        la Martingala ofrece ganancias frecuentes y pequeñas, con el riesgo
        de pérdidas catastróficas en rachas adversas. La Paroli invierte
        este perfil: pérdidas pequeñas frecuentes y ganancias grandes
        ocasionales. D'Alembert y Fibonacci se sitúan en posiciones
        intermedias.

  \item \textbf{El capital finito es determinante:} bajo capital infinito
        (ideal), las diferencias entre estrategias son principalmente de
        volatilidad. Bajo capital finito (real), la Martingala genera
        bancarrotas con mucha mayor frecuencia que la Paroli, lo que en
        la práctica hace que el jugador pierda todo su capital antes de
        poder recuperarse.

  \item \textbf{La frsa converge independientemente de la estrategia:}
        la frecuencia relativa de victorias tiende a $18/37$ para todas las
        estrategias, ya que la probabilidad de cada tirada no depende de las
        apuestas realizadas.

  \item \textbf{Implicación práctica:} apostar en la ruleta es,
        matemáticamente, un juego de suma negativa para el jugador.
        La simulación corrobora cuantitativamente este resultado y permite
        visualizar con claridad por qué ninguna estrategia de apuesta
        puede generar ganancias esperadas positivas a largo plazo.
\end{enumerate}

% ---------------------------------------------------------------------------
\bibliographystyle{unsrt}
\bibliography{referencias}

\end{document}
